\chapter{Finite random-clock certificates for classical realization size}
\label{ch:finite-erlang}

An observed damped oscillation can require a large classical realization even
when a qubit produces it. This fact and the relevant spectral obstruction are
established \cite{rrWolf2009}. The practical difficulty is different: finite,
noisy observations do not determine an exact pole. We therefore measure a finite
list of bounded probabilities whose distance from every small classical model
can be certified directly. The resulting tolerance is conservative; at high
dimension it is much too small for a practical experiment.

\section{The operational contract}
Let $Q$ be a conservative, homogeneous generator with at most $n$ hidden states,
$\rho$ an unknown but fixed initial probability row, and $f\in[0,1]^n$ an unknown
but fixed binary response vector. The detector may be probabilistic. Neither
stationarity, reversibility nor a bound on the exit rates is assumed. Each trial
starts from $\rho$ and measures $f$ at a time chosen independently of the physical
process. The clock does not change the generator, preparation or detector.

Fix a rate $s>0$. For $k\geq1$, choose an Erlang clock $T_k$ with shape $k$ and
rate $s$; set $T_0=0$. Resolvent integration gives
\begin{equation}
 m_k=\mathbb E[\rho e^{QT_k}f]=\rho R^k f,
 \qquad R=s(sI-Q)^{-1}.
 \label{eq:erlang-resolvent}
\end{equation}
The matrix $R$ is stochastic. Thus independently randomized observation times
convert a continuous-time record into finitely many powers of one stochastic
matrix, without estimating derivatives or the hidden exit rates.

The comparison signal is the $+Z$ occupation probability of a qubit prepared
in $+Z$, rotating under $\omega X/2$ with isotropic depolarization rate $\gamma$:
\begin{equation}
 p_*(t)=\frac{1+e^{-\gamma t}\cos(\omega t)}2,
 \qquad
 m_k^*=\frac{1+\operatorname{Re}\mu^k}2,
 \qquad \mu=\frac{s}{s+\gamma-i\omega}=u+iv.
 \label{eq:erlang-qubit}
\end{equation}
Here $\omega,s>0$, $\gamma\geq0$, and $v>0$. Rational rates give rational
probabilities $m_k^*$, even though the observation times are continuous.

\section{A finite robustness theorem}
The established generator spectral bound is
\begin{equation}
 |\operatorname{Im}\lambda|
 \leq-\operatorname{Re}\lambda\cot(\pi/n).
 \label{eq:generator-sector-rf}
\end{equation}
Writing $t_n=\tan(\pi/n)$ shows that every eigenvalue $z$ of $R$ lies in the lens
\begin{equation}
 |z|^2-\operatorname{Re}z+t_n|\operatorname{Im}z|\leq0.
 \label{eq:resolvent-lens}
\end{equation}
In particular, for any certified $0<t_0\leq t_n$, the lens is contained
in both reflected disks with centers $1/2\pm it_0/2$ and radius
$\sqrt{1+t_0^2}/2$.

For the sharper bound it is useful to normalize the characteristic polynomial.
Define
\[
 \nu=\frac{1-\mu}{1+\mu}=a-ib,\qquad b>0,\qquad
 K=\sqrt{1+(1+u)^2/v^2},\qquad q=\lfloor(n-1)/2\rfloor.
\]
If $\gamma/\omega<t_0$, then $\nu$ is outside the upper-centered bounding disk.
Choose
\[
 0<d\leq\min\left\{b,
 \left|\nu-\left(\tfrac12+\tfrac{it_0}{2}\right)\right|
 -\frac{\sqrt{1+t_0^2}}2\right\},\qquad
 r_0=\frac{v}{|1+\mu|},\qquad
 h=\frac{|1+\mu|^2}{4}d(2b-d).
\]
Notice that $h\leq r_0^2$.

\begin{theorem}[Finite Erlang exclusion]
\label{thm:finite-erlang}
Under the preceding assumptions, every classical model with at most $n\geq3$
states satisfies
\begin{equation}
 \max_{0\leq k\leq n+1}|m_k-m_k^*|
 \geq\Delta_n,
 \qquad
 \Delta_n=\frac{|1-\mu|\,r_0^{n-1-2q}h^q}{4K}>0.
 \label{eq:finite-erlang-gap}
\end{equation}
Only $n+2$ observation settings enter this inequality. No fitted pole, exact
rank determination, spectral visibility assumption, or exit-rate bound is used.
\end{theorem}

\begin{proof}
Pad a smaller model by unreachable absorbing states. Let $P$ be the real monic
characteristic polynomial of its $n\times n$ stochastic resolvent $R$, and let
$C$ be the sum of the absolute values of its coefficients. The resolvent roots
have positive real parts. Each real-linear and conjugate-quadratic factor has
alternating coefficients, and products preserve those signs. Hence
$C=|P(-1)|$. Also $P(1)=0$ and $P(R)=0$.

If the maximum moment error is $\varepsilon$, Cayley--Hamilton applied at shifts
zero and one gives
\[
 |\operatorname{Re}P(\mu)|\leq2\varepsilon C,
 \qquad
 |\operatorname{Re}(\mu P(\mu))|\leq2\varepsilon C.
\]
The constant half of~\eqref{eq:erlang-qubit} cancels because $P(1)=0$.
Writing $P(\mu)=A+iB$ yields
$|B|\leq2\varepsilon C(1+u)/v$ and therefore
$|P(\mu)|/C\leq2\varepsilon K$.

For any root $z=s/(s-\lambda)$, put
\[
 w=\frac{1-z}{1+z}=1-\frac{2s}{2s-\lambda}.
\]
The lens~\eqref{eq:resolvent-lens} is invariant under $z\mapsto1-z$, so $w$ lies
in the same lens. Direct algebra gives
\[
 \frac{|\mu-z|}{|1+z|}=\frac{|1+\mu|}{2}|\nu-w|.
\]
A real root therefore contributes at least $r_0$. For a conjugate pair write
$y=|\operatorname{Im}w|$. If $y\geq b-d$, the near factor is at least $d$ and
the far factor at least $b+y\geq2b-d$. If $y\leq b-d$, their product is at least
$(b-y)(b+y)\geq d(2b-d)$. Thus each normalized conjugate pair contributes at
least $h$. One root $z=1$ contributes $|1-\mu|/2$. Because $h\leq r_0^2$,
the most conservative product uses the maximum number $q$ of pairs:
\[
 \frac{|P(\mu)|}{|P(-1)|}
 \geq\frac{|1-\mu|}{2}r_0^{n-1-2q}h^q.
\]
Combining the two inequalities proves the result. Defective matrices cause no
problem because the argument uses the characteristic polynomial, not an
eigenvector decomposition.
\end{proof}

The distance calculation is entirely reproducible with rational interval
arithmetic. If $F=|\nu-c|^2-r^2>0$, then
\[
 |\nu-c|-r=\frac{F}{|\nu-c|+r}.
\]
Rational upper bounds on the denominator give a lower bound on the distance.
The implementation obtains $t_0$ from Machin's formula for $\pi$, alternating
arctangent bounds and $\tan x\geq x+x^3/3$; for $n=4$ it uses $t_0=1$ exactly.
Integer square roots give directed dyadic bounds for the remaining radicals.

\section{Finite duration, calibration and sampling}
The ideal Erlang clock is unbounded. A laboratory implementation can return an
artificial zero when $T_k>L$, without waiting beyond $L$. Its moment changes by
at most
\begin{equation}
 \tau_k(L)=\Pr(T_k>L)
 =e^{-sL}\sum_{j=0}^{k-1}\frac{(sL)^j}{j!}.
 \label{eq:erlang-tail}
\end{equation}
This tail is bounded above by replacing $e^{-sL}$ with the reciprocal of a
positive partial sum for $e^{sL}$. No hidden dynamical rate appears in the bound.

Let $\widehat m_k$ be sample means from $N$ independent trials per setting,
$r_N$ a simultaneous confidence radius, $b$ a calibrated probability error, and
$\tau\geq\max_k\tau_k(L)$. A valid rejection rule is
\begin{equation}
 \max_k|\widehat m_k-m_k^*|+r_N+b+\tau<\Delta_n.
 \label{eq:erlang-statistical-test}
\end{equation}
Hoeffding and a union bound give confidence at least
$1-2(n+2)e^{-2Nr_N^2}$. The same data cannot first select an unrestricted target
and then use a confidence guarantee proved only for a predeclared target.

The supplied design takes $r_N=\Delta_n/4$ and $\tau\leq\Delta_n/8$. With $b=0$,
both false rejection of the classical class and missed detection of the nominal
qubit target are at most the declared $\alpha$. Under the target, the confidence
event gives total distance at most $2r_N+2\tau<\Delta_n$. For an additional
calibration radius, the corresponding power guarantee requires
$2r_N+2\tau+2b<\Delta_n$.

\section{What the certificate achieves, and its cost}
For $\gamma=1/10$, $\omega=2$ and $s=13/10$, exact arithmetic certifies
\[
 n=4:\qquad\Delta_4>0.0016816962795.
\]
There are six clock settings. A fully bounded design with duration at most
$320/13$ and $\alpha=0.05$ uses a conservative total of $135{,}780{,}096$
records. This large count is a limitation of the current sufficient bound, not
a claim of practical superiority. The result closes a qualitative gap: a
strictly positive finite-data exclusion exists without assuming an exact pole
or bounding the rates of invisible modes.

\begin{center}
\begin{tabular}{r r l}
Excluded state count $n$ & Clock settings & Certified uniform gap\\\hline
4 & 6 & $1.6816962795\times10^{-3}$\\
5 & 7 & $3.2905352735\times10^{-4}$\\
6 & 8 & $8.2677897333\times10^{-5}$\\
8 & 10 & $2.7407114770\times10^{-6}$\\
16 & 18 & $1.5863939618\times10^{-13}$\\
62 & 64 & $1.1657896420\times10^{-117}$
\end{tabular}
\end{center}
All entries concern the same damped qubit; the Erlang rate is chosen separately
for each declared $n$ and recorded in the reproducibility data. The last entries
are constructive mathematical bounds, not recommended experimental plans.

A four-state directed cycle of rate $2$, with fixed response probabilities
$(1,1/2,0,1/2)$, exactly produces
$(1+e^{-2t}\cos2t)/2$. The certificate correctly stops separating at
$\gamma=\omega$ for $n=4$. Larger cyclic classical realizations exist for every
strictly damped sinusoid when sufficiently many states are allowed. Thus the
witness bounds classical realization size; it does not exclude every classical
explanation or convert a state count into a thermodynamic energy cost.

The spectral obstruction, resolvent representation and linear recurrence
machinery are established ingredients. The contribution investigated here is
the explicit finite-setting, bounded-duration robustness certificate and its
checkable error accounting. Historical priority remains open. Improving its
sample complexity is a separate research problem; the finite probability-table
witness elsewhere in this reference supplies a substantially stronger practical
comparison in its effective regime.
